% Physical Review D submission source; compiled September 10, 2026.
% Final author sign-off is required before submission; see package checklist.
\documentclass[aps,prd,reprint,amsmath,amssymb,nofootinbib,
longbibliography]{revtex4-2}
\usepackage[T1]{fontenc}
\usepackage[utf8]{inputenc}
\usepackage{amsthm,bm,booktabs,listings}
\usepackage[hidelinks]{hyperref}
\newtheorem{theorem}{Theorem}
\newtheorem{proposition}[theorem]{Proposition}
\newtheorem{corollary}[theorem]{Corollary}
\newtheorem{lemma}[theorem]{Lemma}
\DeclareMathOperator{\rank}{rank}
\DeclareMathOperator{\Tr}{Tr}
\DeclareMathOperator{\diag}{diag}
\DeclareMathOperator{\spanof}{span}
\newcommand{\HH}{\mathcal H}
\newcommand{\dd}{\mathrm d}
\newcommand{\ee}{\varepsilon}
\newcommand{\norm}[1]{\left\lVert#1\right\rVert}
\newcommand{\ip}[2]{\left\langle#1,#2\right\rangle}
\lstset{basicstyle=\ttfamily\scriptsize,breaklines=true,columns=fullflexible,
  keepspaces=true,showstringspaces=false,language=Python}
\begin{document}
\title{Field reconstruction and curvature identifiability from spherical maps
in a closed spatial geometry}
\author{Andrei Preece}
\email{andre.preece@web.de}
\affiliation{Independent Researcher, Germany}
\author{Boris Batenin}
\email{batenin@web.de}
\affiliation{Independent Researcher, Germany}
\date{September 10, 2026}

\begin{abstract}
An angular map, a resolved spatial volume, and a statistical model of a field
contain different information about a closed universe. We distinguish these
inverse problems for scalar fields with the exact harmonic cutoff
$0\leq n\leq L$ on a round three-sphere. For $q$ concentric spherical maps,
we determine the null space explicitly as a product of $q$ radial factors
times the band of degree $L-q$. Its dimension is $N_{L-q}$, where
$N_j=(j+1)(j+2)(2j+3)/6$ and $N_j=0$ for $j<0$.
The resulting rank bound also applies to any $q$ linear scalar channels
equivariant under rotations about the observer. A resolved layer of angular
half-width $\ee$ is injective, but its squared singular values separate into
the powers $\ee^{2k}$ with multiplicities $(L-k+1)^2$; we give the exact
leading coefficient of the weakest mode. We then prove that a single map
has zero local Fisher information about its angular radius after unrestricted
positive isotropic band powers are profiled out. In contrast, two shells with
independently known radial distances determine the curvature radius from
their highest angular sector, provided that sector is nonzero and the
spectral cutoff is exact. We derive its noise information and a contamination
bound. These are conditional results for specified observation operators,
not a detection of spatial curvature or a derivation of cosmic acceleration.
\end{abstract}
\maketitle

\section{Introduction}
\label{sec:intro}

The geometry of a compact spatial slice does not specify which aspects of
that geometry are identifiable from the observations available to one
observer. In particular, access to a two-dimensional sky is not access to
an open three-dimensional region. This distinction is essential when a
geometrical cosmology is formulated through mode counting, analytic
continuation, or a proposed correspondence between a global field and an
observable domain.

Our earlier conditional reconstruction manuscripts \cite{P09,P10}
distinguish an assumed round $S^3$ geometry from a measured
topology and propose restriction and transfer operators as the next step.
A companion working manuscript \cite{CapDraft} treats finite-band scalar
fields on three-dimensional caps. Here we address the different problem
of observation on spherical surfaces, and distinguish recovery of a field
at fixed geometry from identification of the geometry itself.
No dark-energy, phase-evolution, or Fibonacci-scale postulate is needed
for the results below.

The ingredients have substantial precedents. Polynomial factorization on
latitudes is an established method of spherical interpolation
\cite{Castell}; Gaussian quadrature supplies exact finite integration rules
\cite{DLMFQuad}; spherical concentration operators quantify the stability
of restricted observations \cite{Simons}. Cosmological inference is known
to depend on assumptions about the primordial spectrum \cite{Kinney},
including in recent topology calculations \cite{Samandar}. Relative
responses to a common random realization can evade a sample-variance
limit under suitable assumptions \cite{Abramo}.

Our contribution is a complete calculation within one explicitly defined
$S^3$ observation model: the lost field subspace, its recovery by resolved
radial information, an exact statistical degeneracy, and a conditional
two-shell curvature reconstruction. The factorization and quadrature
arguments are applications of established mathematics. We do not claim
that their underlying principles, spectral discreteness, or cancellation
of common realization amplitudes are newly discovered physical laws.
The precise combined results are offered for independent mathematical
and physical assessment; a search of related literature is not proof of
historical priority.

\section{Geometry, field space, and observation model}
\label{sec:model}

On a fixed round spatial slice of radius $R>0$, use observer-centered
coordinates
\begin{equation}
 \dd s^2=R^2(\dd\theta^2+\sin^2\theta\,\dd\Omega^2),
 \qquad 0\leq\theta\leq\pi .
\end{equation}
Write $\dd\sigma=\dd\Omega/(4\pi)$ and normalize the spatial measure to one:
\begin{equation}
 \dd\mu=\frac{\dd V}{2\pi^2R^3}
 =\frac{2}{\pi}\sin^2\theta\,\dd\theta\,\dd\sigma .
 \label{eq:measure}
\end{equation}
Let $Y_{\ell m}$ be orthonormal in $L^2(S^2,\dd\sigma)$; thus $Y_{00}=1$.
Using real spherical harmonics avoids reality constraints on coefficients.
The complex convention gives the same dimensions and spectra.

Let $\HH_n$ be the scalar eigenspace of $-\Delta$ with eigenvalue
$n(n+2)/R^2$. Its dimension is $(n+1)^2$. Our field space is the
\emph{complete} band
\begin{align}
 \HH_{\leq L}&=\bigoplus_{n=0}^L\HH_n,\\
 N_L&=\sum_{n=0}^L(n+1)^2
 =\frac{(L+1)(L+2)(2L+3)}6 .
 \label{eq:dimension}
\end{align}
Set $\HH_{\leq j}=\{0\}$ and $N_j=0$ for $j<0$.
The cutoff is an assumption about the field, not a consequence of
compactness. Compactness alone leaves infinitely many eigenmodes.

An orthonormal basis is
\begin{align}
 Q_{n\ell m}(\theta,\Omega)&=r_{n\ell}(\theta)Y_{\ell m}(\Omega),\\
 r_{n\ell}(\theta)&=a_{n\ell}(\sin\theta)^\ell
 C_{n-\ell}^{\ell+1}(\cos\theta),\\
 a_{n\ell}^2&=
 \frac{2^{2\ell}(n-\ell)!(n+1)(\ell!)^2}{(n+\ell+1)!},
 \quad 0\leq\ell\leq n\leq L.
 \label{eq:radial}
\end{align}
The normalization follows by Gegenbauer orthogonality with weight
$(1-t^2)^{\ell+1/2}$ \cite{DLMFPoly}. In particular,
$r_{n0}=C_n^1$ and $r_{LL}=a_{LL}\sin^L\theta$.

Since Gegenbauer polynomials of degrees zero through $L-\ell$ form a
polynomial basis, every field has the unique representation
\begin{equation}
 u=\sum_{\ell=0}^L\sum_{m=-\ell}^{\ell}
 (\sin\theta)^\ell p_{\ell m}(\cos\theta)Y_{\ell m},
 \qquad \deg p_{\ell m}\leq L-\ell .
 \label{eq:poly}
\end{equation}
This elementary separation is the basis of the proofs.

All surface maps below are scalar values of the \emph{same} field on one
spatial slice, with complete angular coverage and a known common center.
Their norm uses $\dd\sigma$, not physical surface area. Real cosmological
maps generally lie on a past light cone and involve source evolution,
line-of-sight integrals, lensing, and selection effects. Those operations
must be included in a forward model before an observational application.

\section{An exact ceiling on information in angular channels}
\label{sec:rank}

Let $T$ be a linear observation operator with $q$ scalar sky maps as its
output. Suppose it commutes with the $SO(3)$ rotations fixing the observer.
In each angular sector, the field contains $L-\ell+1$ independent radial
coefficients, whereas the output contains at most $q$.

\begin{theorem}[Angular-channel rank bound]
\label{thm:ceiling}
For an operator with these assumptions,
\begin{align}
 \rank T&\leq\sum_{\ell=0}^L(2\ell+1)
       \min\{q,L-\ell+1\}\\
 &=N_L-N_{L-q} .
 \label{eq:ceiling}
\end{align}
Consequently at least $N_{L-q}$ field directions are unidentifiable.
\end{theorem}

\begin{proof}
The restriction of the field representation to rotations about the observer
is
\begin{equation}
 \HH_{\leq L}\cong\bigoplus_{\ell=0}^L
 \left(\mathbb R^{L-\ell+1}\otimes V_\ell\right),
 \qquad \dim V_\ell=2\ell+1.
\end{equation}
An equivariant linear map acts in this sector by a radial matrix
$T_\ell:\mathbb R^{L-\ell+1}\to\mathbb R^q$, tensor the identity on
$V_\ell$. The rank of this matrix is at most the smaller dimension.
For $q\leq L$, the deficit is
$\sum_{\ell=0}^{L-q}(2\ell+1)(L-q-\ell+1)=N_{L-q}$.
For $q\geq L+1$ the deficit is zero.
\end{proof}

This statement includes finitely many radially averaged maps and any
known rotationally invariant transfer functions. Widening a radial
window does not increase the number of output channels. Increasing the
angular pixel count beyond that required to represent the band cannot
remove this radial rank deficit. A sky mask or an additional deterministic
compression of such outputs cannot increase their rank.

\subsection{Concentric shells attain the ceiling}

Choose distinct $\chi_j\in(0,\pi)$, $1\leq j\leq q$, and define
\begin{equation}
 B_q u=\{u(\chi_j,\cdot)\}_{j=1}^q,
 \qquad P_q(t)=\prod_{j=1}^q(t-\cos\chi_j).
\end{equation}
Arbitrary fixed positive weights may be used for the output maps; the
kernel does not depend on them.

\begin{theorem}[Complete shell null space]
\label{thm:kernel}
The observation kernel is exactly
\begin{equation}
 \boxed{\ker B_q=P_q(\cos\theta)\,\HH_{\leq L-q}.}
 \label{eq:kernel}
\end{equation}
Its dimension is $N_{L-q}$ and its rank is $N_L-N_{L-q}$.
Thus $L+1$ distinct nonpolar shells are necessary and sufficient to
recover every field in the complete band.
\end{theorem}

\begin{proof}
Project a zero shell map onto $Y_{\ell m}$. Since $\sin\chi_j\ne0$,
Eq.~\eqref{eq:poly} gives $p_{\ell m}(\cos\chi_j)=0$ for all $j$.
Distinct roots imply that $p_{\ell m}$ is divisible by $P_q$, or is zero
when its degree is less than $q$. The quotient has degree at most
$L-q-\ell$, so the collection of quotients is precisely the polynomial
representation of $\HH_{\leq L-q}$. Conversely every product in
Eq.~\eqref{eq:kernel} belongs to $\HH_{\leq L}$ and vanishes on all the
shells. Multiplication by $P_q(\cos\theta)$ is injective on the quotient
space: a product zero everywhere has its second factor zero on the open
set where $P_q\ne0$, hence zero as a finite harmonic sum. The dimension
and rank follow.
\end{proof}

For a single shell, the rank is only $(L+1)^2$, compared with
$N_L\sim L^3/3$ for the spatial field. The constant mode is included in
these counts. Deleting a measured monopole or dipole gives a different
operator with additional missing information.

The necessity of $L+1$ shells is stronger than a comparison of total
input and output dimensions: the monopole sector alone contains $L+1$
radial degrees of freedom. Angular modes cannot replace those measurements
when the operator preserves observer-centered rotations.

\begin{corollary}[Normal-derivative data]
\label{cor:jets}
If at shell $j$ one measures the derivatives of orders
$0,\ldots,s_j-1$ in the normal direction and $q=\sum_j s_j$, then
\begin{equation}
 \ker B_{\rm jet}=
 \prod_j(\cos\theta-\cos\chi_j)^{s_j}\HH_{\leq L-q}.
 \label{eq:jets}
\end{equation}
\end{corollary}

\begin{proof}
At a nonpolar radius $\theta\mapsto\cos\theta$ is a local coordinate
change with nonzero derivative. The factor $\sin^\ell\theta$ is smooth
and nonzero. Multiplication by that factor and the coordinate change
preserve the order of vanishing. A zero normal jet therefore imposes
the indicated root multiplicity on each $p_{\ell m}$. Apply polynomial
division as before. The physical normal derivative $R^{-1}\partial_\theta$
changes nonzero scaling factors, not the kernel.
\end{proof}

\subsection{An exact, optimally conditioned shell design}

\begin{proposition}[Gaussian shell isometry]
\label{prop:isometry}
Take $q=L+1$ shells with
\begin{equation}
 \chi_j=\frac{j\pi}{L+2},\qquad
 w_j=\frac{2}{L+2}\sin^2\chi_j,
 \quad j=1,\ldots,L+1.
 \label{eq:design}
\end{equation}
Then for every $u\in\HH_{\leq L}$,
\begin{equation}
 \sum_j w_j\int_{S^2}|u(\chi_j,\Omega)|^2\dd\sigma
 =\norm{u}_{L^2(\mu)}^2 .
 \label{eq:isometry}
\end{equation}
Every singular value of this weighted observation operator is one.
\end{proposition}

\begin{proof}
Angular orthogonality reduces the integrand to
$\sum_{\ell m}(1-t^2)^\ell|p_{\ell m}(t)|^2$, a real polynomial of degree
at most $2L$. Equation~\eqref{eq:measure} induces the normalized weight
$(2/\pi)\sqrt{1-t^2}\,\dd t$. The nodes and weights in
Eq.~\eqref{eq:design} are its $(L+1)$-point Gauss--Chebyshev rule of the
second kind, exact through degree $2L+1$ \cite{DLMFQuad}. This proves
the norm identity. Polarization gives the inner-product identity.
\end{proof}

This is an application of Gaussian quadrature, not a new quadrature rule.
Its significance here is that the minimum shell count and perfect
conditioning can be attained simultaneously if the full radial range is
available and the geometry is known. These shells are not claimed to be
accessible cosmological survey surfaces.

\section{A resolved thin layer: uniqueness and its cost}
\label{sec:thin}

Fix $\chi_0\in(0,\pi)$ and let
$0<\ee<\min\{\chi_0,\pi-\chi_0\}$. In this section the data resolve
the radial coordinate throughout the layer, rather than averaging it
into one map. Let $w\geq0$ be an integrable weight on $[-1,1]$, with
$\int w\dd s=1$, positive almost everywhere on some interval. Define
\begin{equation}
 \norm{A_\ee u}^2=
 \int_{-1}^1 w(s)\int_{S^2}
 |u(\chi_0+\ee s,\Omega)|^2\dd\sigma\,\dd s.
 \label{eq:layernorm}
\end{equation}
Use the full spatial norm \eqref{eq:measure} on the input.
All following asymptotics hold at fixed $L$, $\chi_0$, and $w$ as
$\ee\to0^+$. They are not uniform large-bandwidth statements.

\begin{theorem}[Complete thin-layer spectrum]
\label{thm:thin}
The operator $A_\ee$ is injective. Its squared singular values within
angular sector $\ell$ can be ordered as
\begin{equation}
 \lambda_{\ell k}(\ee)=\Theta(\ee^{2k}),
 \quad k=0,\ldots,L-\ell,
 \label{eq:powers}
\end{equation}
each repeated $2\ell+1$ times. In the complete spectrum, the scaling
power $\ee^{2k}$ has multiplicity $(L-k+1)^2$.
For $L\geq1$, the unique weakest branch has $\ell=0$, $k=L$ and
\begin{equation}
 \lambda_{\min}(\ee)\sim
 \frac{\sin^{2L}\chi_0\,h_L(w)}{J_L(\cos\chi_0)}\,\ee^{2L},
 \label{eq:mincoefficient}
\end{equation}
where
\begin{align}
 h_L(w)&=\min_{b_0,\ldots,b_{L-1}}
 \int_{-1}^1w(s)\left|s^L+\sum_{j<L}b_js^j\right|^2\dd s,\\
 J_L(c)&=\frac2\pi\int_{-1}^1(t-c)^{2L}\sqrt{1-t^2}\,\dd t.
 \label{eq:hJ}
\end{align}
For $L=0$, $\lambda_{\min}=1$ and the same coefficient formula applies.
\end{theorem}

\begin{proof}
Use in sector $\ell$ the basis
$\sin^\ell\theta(\cos\theta-c)^i$, $0\leq i\leq M=L-\ell$,
where $c=\cos\chi_0$. Denote its full spatial Gram matrix by $F_\ell$.
It is fixed and positive definite. Its observation Gram matrix obeys
\begin{equation}
 G_\ell(\ee)=D_\ee M_\ell(\ee)D_\ee,
 \quad D_\ee=\diag(1,\ee,\ldots,\ee^M).
 \label{eq:gramfactor}
\end{equation}
Indeed, uniform Taylor expansion on $s\in[-1,1]$ gives
\begin{equation}
 \begin{aligned}
 \bigl[M_\ell(0)\bigr]_{ij}
 &=\sin^{2\ell}\chi_0(-\sin\chi_0)^{i+j}\mu_{i+j},\\
 \mu_r&=\int_{-1}^1s^rw(s)\dd s .
 \end{aligned}
 \label{eq:momentlimit}
\end{equation}
The limiting moment matrix is positive definite because a nonzero
polynomial cannot vanish on an interval. Continuity therefore gives
fixed positive upper and lower matrix bounds on $M_\ell(\ee)$ for
sufficiently small $\ee$. Comparing the Rayleigh quotient
$a^*G_\ell a/(a^*F_\ell a)$ with $a^*D_\ee^2a/(a^*a)$ and applying
the finite-dimensional min--max principle proves Eq.~\eqref{eq:powers}.
At each fixed admissible $\ee$, the same polynomial argument proves
positive definiteness and injectivity. Summing $2\ell+1$ over
$0\leq\ell\leq L-k$ gives $(L-k+1)^2$.

For the weakest branch only $\ell=0$ contributes at exponent $2L$.
The reciprocal eigenvalues are those of
$F_0^{1/2}G_0^{-1}F_0^{1/2}$. From Eq.~\eqref{eq:gramfactor},
\begin{equation}
 \ee^{2L}G_0^{-1}\longrightarrow
 [M_0(0)^{-1}]_{LL}e_Le_L^*.
\end{equation}
Its nonzero eigenvalue after multiplication by $F_0^{1/2}$ on both sides
is $[M_0(0)^{-1}]_{LL}[F_0]_{LL}$.
The Schur complement formula for the monic orthogonal polynomial gives
$( [\mu_{i+j}]^{-1})_{LL}=1/h_L(w)$. Equation~\eqref{eq:momentlimit}
then yields
$[M_0(0)^{-1}]_{LL}=\sin^{-2L}\chi_0/h_L(w)$, while
$[F_0]_{LL}=J_L(c)$. Taking the reciprocal proves
Eq.~\eqref{eq:mincoefficient}.
\end{proof}

For the uniform weight $w=1/2$, the leading coefficient of the Legendre
polynomial gives
\begin{equation}
 h_L(1/2)=\frac{2^{2L}}{(2L+1)\binom{2L}{L}^2}.
 \label{eq:legendre}
\end{equation}
The spatial moment can be evaluated without quadrature:
\begin{equation}
 J_L(c)=\sum_{r=0}^L\binom{2L}{2r}c^{2L-2r}
 \frac{\binom{2r}{r}}{4^r(r+1)}.
 \label{eq:catalan}
\end{equation}
At the equator and $L=1$, for example,
$\lambda_{\min}\sim4\ee^2/3$. The exponents describe multiplicities of
asymptotic orders, not exact degeneracies between different $\ell$.

The inverse norm grows as
\begin{equation}
 \norm{A_\ee^+}\sim
 \left[\frac{J_L(\cos\chi_0)}{\sin^{2L}\chi_0h_L(w)}\right]^{1/2}
 \ee^{-L}.
 \label{eq:inversecost}
\end{equation}
Here the pseudoinverse acts on the finite-dimensional signal range.
Thus the infinite precision needed in the zero-width limit is quantified,
even though any positive resolved thickness gives uniqueness.
In the limit, exactly $(L+1)^2$ singular-value branches remain of order
one, matching the rank of one shell.

The normalization in Eq.~\eqref{eq:layernorm} matters. For the actual
unrenormalized spatial norm restricted to a uniform thin layer,
\begin{equation}
 \int_{\chi_0-\ee}^{\chi_0+\ee}\frac2\pi\sin^2\theta\dd\theta
 \int|u|^2\dd\sigma,
\end{equation}
every exponent increases by one. The weakest coefficient is multiplied
by $4\sin^2\chi_0/\pi$ relative to Eq.~\eqref{eq:mincoefficient} with
$w=1/2$. A physically specified noise model must determine which norm
describes the experiment.

\subsection{Noise and indistinguishability}
\label{sec:noise}

For finitely many real signal coefficients with Gaussian additive noise
$\mathcal N$ of positive-definite covariance, the Fisher matrix is
$F=T^*\mathcal N^{-1}T$. Hence $\ker F=\ker T$. If $v\in\ker T$ has
unit norm, the fields $u+av$ and $u-av$ have identical data distributions.
For any estimator, their maximum mean squared error is at least $a^2$:
the average of their two risks equals the risk about their midpoint plus
$a^2$. No assertion about exact reconstruction can overcome this null
space without additional information.

For resolved-layer data, project the observations onto an orthonormal
basis of the finite-dimensional signal range and assume coefficient
noise of variance $\sigma^2$. This specifies the otherwise ambiguous
notion of white continuum noise. The Cram\'er--Rao variance in the weakest
direction is $\sigma^2/\lambda_{\min}$ and therefore grows as
$\ee^{-2L}$. A different noise covariance uniformly bounded above and
below, independently of $\ee$, preserves the powers but changes the
coefficients.

\section{A single map cannot identify its radius with free band powers}
\label{sec:degeneracy}

Now let the input be a zero-mean $SO(4)$-isotropic real Gaussian field.
In the orthonormal basis \eqref{eq:radial}, its coefficients satisfy
\begin{equation}
 \mathbb E[x_{n\ell m}x_{n'\ell'm'}]
 =P_n\delta_{nn'}\delta_{\ell\ell'}\delta_{mm'},
 \qquad P_n>0.
 \label{eq:prior}
\end{equation}
There is one unrestricted band power per spatial degree $n$.
For a shell at angular radius $\chi$, the sky covariance is diagonal:
\begin{align}
 C_\ell(\chi,P)&=\sum_{n=\ell}^L M_{\ell n}(\chi)P_n,\\
 M_{\ell n}(\chi)&=r_{n\ell}(\chi)^2,\quad n\geq\ell,
 \label{eq:mixing}
\end{align}
and $M_{\ell n}=0$ for $n<\ell$.

\begin{theorem}[Exact local statistical degeneracy]
\label{thm:degeneracy}
Fix $\chi_0\in(0,\pi)$ and any strictly positive vector $P^{(0)}$.
There is an open interval around $\chi_0$ on which
\begin{equation}
 P(\chi)=M(\chi)^{-1}M(\chi_0)P^{(0)}>0
 \label{eq:degcurve}
\end{equation}
and the entire Gaussian sky distribution is unchanged.
The Fisher information for $\chi$ after profiling over all $P_n$ is
identically zero at every such interior parameter point.
If an independently fixed radial distance $d$ defines $\chi=d/R$,
the same local nonidentifiability applies to $R$.
\end{theorem}

\begin{proof}
The matrix $M$ is upper triangular with diagonal
\begin{equation}
 M_{\ell\ell}=a_{\ell\ell}^2\sin^{2\ell}\chi>0.
\end{equation}
It is therefore invertible and smooth for all $0<\chi<\pi$.
Equation~\eqref{eq:degcurve} preserves $C=MP$ exactly. Positivity is
preserved in a sufficiently small interval by continuity and the finite
number of strictly positive components. A zero-mean Gaussian field is
determined by its covariance, so this equality concerns the full map
likelihood, not just a single realization or an approximate fit.

For the Fisher statement, the derivative $M'(\chi)P$ lies in the span
of the band-power derivative columns of $M$. Thus the score for $\chi$
is a linear combination of the nuisance-parameter scores. Its residual
after projection onto their span vanishes. Equivalently,
$P'=-M^{-1}M'P$ is a tangent vector of constant likelihood in
Eq.~\eqref{eq:degcurve}. The change of variable $\chi=d/R$ preserves
the degeneracy.
\end{proof}

This result is local in the interior of the positive-power cone. It does
not assert that every positive sky spectrum is realizable at every radius.
At a boundary where some $P_n=0$, positivity can constrain a variation.
A restricted primordial spectrum, dynamics, cross-correlated channels,
or external distance information can change the problem. Known additive
noise independent of $\chi$ does not remove the equality of distributions.
A parameter-independent mask or other deterministic processing also
preserves it.

There is a further exact ambiguity even with fixed band powers:
$M(\pi-\chi)=M(\chi)$. This follows from the parity of Gegenbauer
polynomials. An angular-radius prior restricting the domain can remove
this antipodal pair, but not the local free-spectrum degeneracy above.

For $L=1$, the model is particularly transparent:
\begin{equation}
 C_0=P_0+4\cos^2\chi\,P_1,
 \qquad C_1=\frac43\sin^2\chi\,P_1.
\end{equation}
Changing $\chi$ can be compensated by
$P_1=3C_1/(4\sin^2\chi)$ and
$P_0=C_0-3C_1\cot^2\chi$, while these remain positive.

The theorem concerns a pure shell-value model on a fixed complete finite
band. It is not a no-go theorem for curvature measurements using the
actual cosmic microwave background (CMB), whose acoustic transfer
functions, polarization, and cosmological priors supply other structure.
It does show why a spectral gap or a projected angular spectrum alone
cannot establish curvature without specifying the allowed source spectra.

\section{Two shells can identify curvature without reconstructing the field}
\label{sec:twoshell}

Suppose two geodesic radii $0<d_1<d_2$ on the same slice are independently
calibrated in common length units. Let
\begin{equation}
 \kappa=\frac{d_2}{d_1}>1,\qquad
 \chi=\frac{d_1}{R},\qquad R>\frac{d_2}{\pi}.
 \label{eq:branch}
\end{equation}
The branch condition places both shells before the antipode. It is
necessary for the geometric interpretation used here.

\begin{theorem}[Conditional two-shell curvature reconstruction]
\label{thm:radius}
Assume the exact cutoff $L\geq1$ and a nonzero projection of the field
onto the angular sector $\ell=L$ about the common observer. Let $b_j$
be the real vector of its $2L+1$ measured coefficients on shell $j$.
Then
\begin{equation}
 b_2=g(R)b_1,\qquad
 g(R)=\left[\frac{\sin(d_2/R)}{\sin(d_1/R)}\right]^L.
 \label{eq:ratio}
\end{equation}
The measured ratio
$g=\ip{b_1}{b_2}/\norm{b_1}^2$ determines exactly one radius in
Eq.~\eqref{eq:branch}, whenever $0<g^{1/L}<\kappa$.
This determination requires neither the unknown realization coefficients
nor the positive band power $P_L$.
\end{theorem}

\begin{proof}
Only $n=L$ contributes to angular degree $\ell=L$ in the complete band.
Equation~\eqref{eq:radial} therefore gives
$b_j=a_{LL}\sin^L(d_j/R)x_L$, where $x_L$ is the common nonzero vector.
This proves Eq.~\eqref{eq:ratio} and the amplitude cancellation.

It remains to prove uniqueness, not just write an inversion equation.
Define $f_\kappa(x)=\sin(\kappa x)/\sin x$ on $0<x<\pi/\kappa$.
Its logarithmic derivative is
\begin{equation}
 \frac{f_\kappa'(x)}{f_\kappa(x)}
 =\kappa\cot(\kappa x)-\cot x<0.
 \label{eq:monotonic}
\end{equation}
To establish the inequality, put $h(x)=x\cot x$.
Then $h'(x)=[\sin x\cos x-x]/\sin^2x<0$ on $(0,\pi)$:
the numerator is zero at zero and has derivative $-2\sin^2x$.
Since $\kappa x>x$, $h(\kappa x)<h(x)$, proving
Eq.~\eqref{eq:monotonic}. Finally,
$f_\kappa(0^+)=\kappa$ and $f_\kappa((\pi/\kappa)^-)=0$.
Every admissible ratio therefore has exactly one $x$, and $R=d_1/x$.
\end{proof}

In the isotropic Gaussian model with $P_L>0$, the required vector is
nonzero with probability one. Deterministic fields can have a zero
projection in that sector even if they contain other degree-$L$ modes;
the theorem explicitly excludes this case.

For $\kappa=2$, there is a closed form:
\begin{equation}
 \boxed{R=\frac{d_1}{\arccos\left(g^{1/L}/2\right)}},
 \qquad 0<g^{1/L}<2.
 \label{eq:closedform}
\end{equation}
The same ratio can be written from exact ensemble covariances as
$g=C_L^{12}/C_L^{11}$, and
$(C_L^{12})^2=C_L^{11}C_L^{22}$. In the ideal noiseless model it already
holds realization by realization, so it does not require an ensemble
estimate. The cancellation uses the same principle of relative response
to a common realization that appears in multitracer inference
\cite{Abramo}; the sine-ratio inversion is specific to the present
highest-sector shell model.

For $L\geq2$, two shells still leave $N_{L-2}$ field directions
undetermined at the recovered radius. Thus identification of one
geometrical parameter is compatible with failure of full field
reconstruction. For $L=1$, two shells also suffice for the field once
the geometry has been determined.

\subsection{Noise information and the nearly flat limit}

Write the two observed top-sector vectors as
\begin{equation}
 y_1=a+z_1,\qquad y_2=g(R)a+z_2,
 \qquad z_j\sim\mathcal N(0,\sigma_j^2 I),
 \label{eq:twonoise}
\end{equation}
with independent noises and an unrestricted deterministic nuisance vector
$a\ne0$. Taking the Schur complement over $a$ gives the information
\begin{equation}
 I_g=\frac{\norm{a}^2}{\sigma_2^2+g^2\sigma_1^2},
 \qquad
 I_R=\frac{[g'(R)]^2\norm{a}^2}{\sigma_2^2+g^2\sigma_1^2}.
 \label{eq:Ifisher}
\end{equation}
For completeness, before profiling the blocks are
$I_{gg}=\norm{a}^2/\sigma_2^2$,
$I_{ga}=g a^T/\sigma_2^2$, and
$I_{aa}=(\sigma_1^{-2}+g^2\sigma_2^{-2})I$.
These produce Eq.~\eqref{eq:Ifisher} directly. This is a conditional
noise calculation, not a forecast for a real survey or a claim about
the bias of the simple ratio estimator in noise.

Let $K=R^{-2}$ denote the positive sectional curvature, distinct from
the cosmological parameter $\Omega_K$. For fixed calibrated distances,
\begin{equation}
 \ln g=L\ln\kappa-\frac{L(d_2^2-d_1^2)}6K+O(K^2).
 \label{eq:flat}
\end{equation}
Thus the ratio tends to its flat-space value $\kappa^L$, and
$\partial_R\ln g=L(d_2^2-d_1^2)/(3R^3)+O(R^{-5})$.
At fixed nonzero reference signal $a$ and noise levels, $I_R$ decreases
as $R^{-6}$. The information for $K$ can remain finite at $K=0$ under
this normalization, but finite precision cannot imply a relative
precision on an arbitrarily small curvature. If the physical source
amplitude changes with $R$, it must also be included in the forecast.

\subsection{What breaks the reconstruction}
\label{sec:leakage}

The exact upper spectral edge is central. If modes with $n>L$ are
present, they contribute to $\ell=L$ with different radial factors and
the common-vector relation need not hold. Applying an angular low-pass
filter does not impose the required three-dimensional spectral cutoff.
Let the two noiseless maps be
\begin{equation}
 y_1=a+e_1,\qquad y_2=g a+e_2,
\end{equation}
where $e_j$ includes any such contamination. Define
$\widehat g=\ip{y_1}{y_2}/\norm{y_1}^2$. Then the exact identity
\begin{equation}
 \widehat g-g=\frac{\ip{y_1}{e_2-ge_1}}{\norm{y_1}^2}
\end{equation}
implies, if $\norm{e_j}\leq\eta_j\norm{a}$ and $\eta_1<1$,
\begin{equation}
 |\widehat g-g|\leq\frac{\eta_2+g\eta_1}{1-\eta_1}.
 \label{eq:leakbound}
\end{equation}
This bound requires externally justified contamination bounds; it cannot
be used to infer them from the two maps alone. On a compact interval of
admissible $R$, monotonicity converts a ratio error into a radius error
using the minimum of $g'(R)>0$ on that interval.

Unknown relative calibration is another exact obstruction: if the maps
have an unknown gain ratio $b$, the observable ratio is $b\,g(R)$ and
one ratio cannot determine both parameters. Source evolution between
two light-cone shells can produce the same problem. Finally, $d_1$ and
$d_2$ must not be obtained by assuming the curvature that is being tested.
Known $\kappa$ alone gives $\chi$; one independently known absolute
distance is needed to turn it into a dimensionful radius.

\section{Relation to cosmological reconstruction}
\label{sec:cosmology}

Theorems~\ref{thm:kernel}, \ref{thm:thin}, \ref{thm:degeneracy}, and
\ref{thm:radius} answer different questions. At fixed geometry, finitely
many maps have an exact field null space. A resolved layer removes that
null space with a calculable loss of stability as its width decreases.
With unknown geometry, a freely adjustable positive spectrum can erase
single-map sensitivity to radius exactly. A second calibrated response
to the same highest-sector realization can restore that sensitivity
without restoring all field coefficients.

No assumed volume fraction is an information fraction in these statements.
For example, the cap fraction
\begin{equation}
 f(\chi)=\frac{\chi-\sin\chi\cos\chi}{\pi}
\end{equation}
does not determine the rank or the noise-conditioned singular values
of a sky-transfer operator. Nor does analytic continuation from an open
volume imply continuation from one embedded two-sphere: the explicit
kernel \eqref{eq:kernel} is a counterexample.

For a physical application, a scalar channel at fixed angular degree
has a transfer matrix of the form
\begin{equation}
 y_{\alpha,\ell m}=\sum_{n=\ell}^L
 T_{\alpha\ell n}(R,\vartheta)x_{n\ell m}+z_{\alpha,\ell m},
 \label{eq:physicalT}
\end{equation}
where $\vartheta$ includes evolution, calibration, and source parameters.
At fixed parameters, Theorem~\ref{thm:ceiling} supplies an exact upper
bound for $q$ such scalar channels. The actual rank may be lower.
For parameter identification, one must additionally differentiate the
likelihood and remove nuisance directions, as in
Theorem~\ref{thm:degeneracy}. An invertible field map at fixed $R$
does not by itself prove that $R$ is identifiable.

The CMB, polarization, and galaxy observations do not have the unmodified
shell transfer used in Theorem~\ref{thm:radius}. Their source fields,
removed low multipoles, and redshift-dependent transfer functions must
be treated explicitly. In particular, neither a sharp cutoff nor its
location has been established by the supplied manuscripts. The present
reconstruction is therefore a conditional benchmark for developing such
an experiment, not an observational determination of the Universe's
radius, topology, dark-energy content, or initial state.

\section{Verification and research methods}
\label{sec:verification}

The analytic arguments are finite-dimensional proofs. Numerical checks
are independent consistency tests of formulas and normalizations, not
substitutes for those proofs. Appendix~\ref{app:code} contains the full
Python program used here, with no input files and no third-party packages.
It uses exact rational row reduction for shell and confluent ranks,
and 160-digit decimal arithmetic for moment matrices, normalization,
singular values, and covariance degeneracies.

The resolved-layer Gram matrices are evaluated by finite trigonometric
moment sums. For uniform radial weight,
\begin{equation}
 \begin{aligned}
 \langle\cos^n\theta\rangle_{\rm layer}
 &=2^{-n}\sum_{j=0}^n\binom nj
 \cos[(n-2j)\chi_0]\\
 &\quad\times\operatorname{sinc}[(n-2j)\ee],\\
 \langle\cos^{2r}\theta\rangle_{\mu}
 &=\frac{\binom{2r}{r}}{4^r(r+1)},\\
 \langle\cos^{2r+1}\theta\rangle_{\mu}&=0 .
 \end{aligned}
 \label{eq:verificationmoments}
\end{equation}
with $\operatorname{sinc}(0)=1$. These expressions avoid a quadrature
resolution error as the layer narrows. Full-norm Cholesky whitening
and symmetric Jacobi diagonalization then give the squared singular
values. No astronomical observations or likelihood chains are used.

The executed program completed 774 assertions. Exact shell ranks were
checked for $0\leq L\leq8$ and $1\leq q\leq L+2$, and confluent ranks
for $0\leq L\leq6$. Gaussian-shell moments through their required degree
agreed with independent semicircle moments to better than $10^{-125}$.
The local positive-spectrum degeneracy was checked for $1\leq L\leq7$,
at $\chi=0.9$ and $0.90001$ radians, with covariance residuals below
$10^{-125}$.

For the thin-layer tests, $\chi_0=1.1$ radians and $w=1/2$.
Slopes from widths $0.04$ and $0.02$ agreed with all predicted powers
through $L=6$ within $0.012$. Table~\ref{tab:thin} reports the weakest
coefficient and its independently computed finite-width ratio.
Two-shell inversions with $d_1=3$, $d_2=6$, and $R=10$ in common length
units recovered $R$ to relative error below $10^{-100}$ for
$L=1,2,4,7$. The Fisher Schur complement was also checked with exact
fractions. These high-precision arithmetic residuals are not error bars
on a physical measurement.

\begin{table}[t]
\caption{Synthetic thin-layer check. $c_L$ is the coefficient in
Eq.~\eqref{eq:mincoefficient}; the last column is
$\lambda_{\min}(0.01)/(c_L\,0.01^{2L})$.}
\label{tab:thin}
\begin{ruledtabular}
\begin{tabular}{ccc}
$L$ & $c_L$ & Finite-width ratio\\
0 & 1 & 1.000000000 \\
1 & 0.58091171 & 0.999970550 \\
2 & 0.11781344 & 1.000037077 \\
3 & 0.018139078 & 1.000061991 \\
4 & 0.0024629588 & 1.000082387 \\
5 & 0.00031074507 & 1.000101175 \\
6 & 3.7342137e-05 & 1.000119207 \\
\end{tabular}
\end{ruledtabular}
\end{table}


\subsection{AI-assisted research disclosure}
OpenAI ChatGPT, accessed through Codex Work Mode on September 9, 2026,
was used substantively for literature synthesis, derivation assistance,
proof checking, and generation and execution of the verification program.
The backend model version was not exposed in the session metadata. Boris
Batenin supplied the source manuscripts, defined the research direction,
and requested explicit proofs and executable checks. Verification of the
AI-assisted outputs in this workflow consists of the analytic derivations
and reproducible numerical consistency tests described in this section;
the packaged program was rerun and completed the 774 assertions reported
above. The named authors retain responsibility for the scientific content
and the final submitted text.

\section{Discussion}

The sharply stated result is an observability distinction. A finite set
of angular channels can have a provable field null space while retaining
information about a geometric parameter. Conversely, a map with infinitely
precise angular coefficients can contain no local information about its
radius when its source spectrum is free. These conclusions follow from
the operator and nuisance space, not from spatial compactness alone.

The strongest route to a physical test is to replace the ideal shell
responses by a calibrated light-cone transfer calculation, assess its
nuisance-projected Fisher information, and bound leakage from modes
outside any effective band. The exact null space, thin-layer coefficient,
Gaussian-shell isometry, and two-shell monotonicity provide separate
checks on that calculation. Failure of a proposed transfer model to pass
one of these checks can reveal an information assumption before it is
mistaken for a cosmological prediction.

\begin{acknowledgments}
OpenAI ChatGPT also assisted with English drafting and LaTeX preparation.
Its substantive research use is disclosed in Sec.~\ref{sec:verification}.
No AI system is listed as an author.
\end{acknowledgments}

\section*{Data Availability}
No observational or experimental data were generated or analyzed in this
study. The synthetic calculations supporting the reported consistency
checks are reproducible from the complete Python verification program in
Appendix~\ref{app:code}. The program requires only the Python standard
library, prints the numerical results, and writes a local JSON report.

\appendix
\section{Complete verification program}
\label{app:code}
Copy the following listing into \texttt{shell\_verification.py} and run
\texttt{python3 shell\_verification.py}. Assertion failures stop the
program. The fixed tolerances test only the finite cases stated in
Sec.~\ref{sec:verification}; the all-$L$ claims are proved in the text.
\begin{lstlisting}
"""Synthetic checks for the spherical-shell manuscript;
no observational input."""
from decimal import Decimal as D, localcontext
from fractions import Fraction as Q
from math import comb, factorial
from pathlib import Path
import json

PI = D('3.14159265358979323846264338327950288419716939937510'
       '58209749445923078164062862089986280348253421170679'
       '82148086513282306647093844609550582231725359408128'
       '4811174502841027019385')


def trig(x, sine=False):
    term = total = x if sine else D(1)
    for r in range(1, 400):
        j = 2*r + int(sine)
        term *= -x*x / D(j*(j-1))
        total += term
        if abs(term) < D('1e-155'):
            return total
    raise ArithmeticError('trigonometric series')


def mul(a, b):
    return [[sum(x*y for x, y in zip(row, col))
             for col in zip(*b)] for row in a]


def trans(a):
    return [list(c) for c in zip(*a)]


def chol(a):
    n = len(a)
    b = [[D(0)]*n for _ in range(n)]
    for i in range(n):
        for j in range(i+1):
            v = a[i][j] - sum(b[i][k]*b[j][k] for k in range(j))
            b[i][j] = v.sqrt() if i == j else v/b[j][j]
    return b


def invlower(a):
    n = len(a)
    b = [[D(0)]*n for _ in range(n)]
    for j in range(n):
        for i in range(n):
            b[i][j] = (D(i == j) - sum(a[i][k]*b[k][j]
                         for k in range(i)))/a[i][i]
    return b


def eig(a):
    a = [r[:] for r in a]
    n = len(a)
    if n == 1:
        return [a[0][0]]
    for _ in range(400*n*n):
        p, q = max(((i, j) for i in range(n) for j in range(i+1, n)),
                   key=lambda ij: abs(a[ij[0]][ij[1]]))
        off = a[p][q]
        if abs(off) < D('1e-125'):
            return sorted((a[i][i] for i in range(n)), reverse=True)
        tau = (a[q][q]-a[p][p])/(2*off)
        t = D(1 if tau >= 0 else -1)/(abs(tau)+(1+tau*tau).sqrt())
        c = 1/(1+t*t).sqrt()
        s = t*c
        for k in range(n):
            if k not in (p, q):
                kp, kq = a[k][p], a[k][q]
                a[k][p] = a[p][k] = c*kp-s*kq
                a[k][q] = a[q][k] = s*kp+c*kq
        a[p][p] -= t*off
        a[q][q] += t*off
        a[p][q] = a[q][p] = D(0)
    raise ArithmeticError('Jacobi convergence')


def N(L):
    return (L+1)*(L+2)*(2*L+3)//6 if L >= 0 else 0


def semicircle(n):
    return (D(comb(n, n//2))/D(n//2+1)/D(2)**n
            if n % 2 == 0 else D(0))


def interval_moment(n, chi, eps):
    ans = D(0)
    for k in range(n+1):
        m = n-2*k
        sinc = trig(m*eps, True)/(m*eps) if m else D(1)
        ans += D(comb(n, k))*trig(m*chi)*sinc
    return ans/D(2)**n


def moment(ell, r, center, base):
    return sum(D(comb(r, j))*(-center)**(r-j)
               * D(comb(ell, a))*(-1)**a*base(j+2*a)
               for j in range(r+1) for a in range(ell+1))


def matrices(L, ell, chi, eps):
    c = trig(chi)
    m = L-ell+1
    full = [[moment(ell, i+j, c, semicircle) for j in range(m)]
            for i in range(m)]
    obs = [[moment(ell, i+j, c,
             lambda n: interval_moment(n, chi, eps)) for j in range(m)]
           for i in range(m)]
    return obs, full


def spectrum(L, ell, chi, eps):
    obs, full = matrices(L, ell, chi, eps)
    b = invlower(chol(full))
    return eig(mul(mul(b, obs), trans(b)))


def rational_rank(a):
    a = [list(map(Q, row)) for row in a]
    row = 0
    for col in range(len(a[0])):
        pivot = next((r for r in range(row, len(a)) if a[r][col]), None)
        if pivot is None:
            continue
        a[row], a[pivot] = a[pivot], a[row]
        v = a[row][col]
        a[row] = [x/v for x in a[row]]
        for r in range(len(a)):
            if r != row:
                v = a[r][col]
                a[r] = [x-v*y for x, y in zip(a[r], a[row])]
        row += 1
        if row == len(a):
            break
    return row


def gegen(k, lam, t):
    if k == 0:
        return D(1)
    old, cur = D(1), 2*lam*t
    for j in range(1, k):
        old, cur = cur, (2*(j+lam)*t*cur-(j+2*lam-1)*old)/(j+1)
    return cur


def radial(n, ell, chi):
    norm2 = (D(2)**(2*ell)*D(factorial(n-ell)*(n+1)*factorial(ell)**2)
             / D(factorial(n+ell+1)))
    return norm2.sqrt()*trig(chi, True)**ell*gegen(n-ell, ell+1, trig(chi))


def mixing(L, chi):
    return [[radial(n, ell, chi)**2 if n >= ell else D(0)
             for n in range(L+1)] for ell in range(L+1)]


def upper_solve(a, b):
    x = [D(0)]*len(b)
    for i in reversed(range(len(b))):
        x[i] = (b[i]-sum(a[i][j]*x[j] for j in range(i+1, len(b))))/a[i][i]
    return x


def verify():
    checks = 0
    thin, rank_rows, curvature_rows = [], [], []
    max_isometry = D(0)
    max_covariance_residual = D(0)
    with localcontext() as ctx:
        ctx.prec = 160
        for L in range(9):
            for q in range(1, L+3):
                nodes = [Q(2*j-q+1, q+1) for j in range(q)]
                rank = sum((2*ell+1)*rational_rank(
                    [[t**k for k in range(L-ell+1)] for t in nodes])
                    for ell in range(L+1))
                assert rank == N(L)-N(L-q)
                checks += 1
            rank_rows.append([L, N(L), N(L)-N(L-1), N(L)-N(L-2)])

            # Gaussian quadrature is checked against independent
            # semicircle moments.
            angles = [D(j)*PI/D(L+2) for j in range(1, L+2)]
            weights = [D(2)/D(L+2)*trig(a, True)**2 for a in angles]
            for ell in range(L+1):
                for r in range(2*(L-ell)+1):
                    observed = sum(w*trig(a, True)**(2*ell)*trig(a)**r
                                   for a, w in zip(angles, weights))
                    full = sum(D(comb(ell, k))*(-1)**k*semicircle(r+2*k)
                               for k in range(ell+1))
                    error = abs(observed-full)
                    max_isometry = max(max_isometry, error)
                    assert error < D('1e-125')
                    checks += 1

        # Confluent ranks, including degrees lower than total jet count.
        for L in range(7):
            for ell in range(L+1):
                rows = []
                for t, multiplicity in [(Q(-1, 3), 2), (Q(2, 5), 3)]:
                    for derivative in range(multiplicity):
                        rows.append([Q(factorial(k), factorial(k-derivative))
                                     * t**(k-derivative)
                                     if k >= derivative else Q(0)
                                     for k in range(L-ell+1)])
                assert rational_rank(rows) == min(5, L-ell+1)
                checks += 1

        for L in range(7):
            counts = {k: 0 for k in range(L+1)}
            largest_slope_error = D(0)
            for ell in range(L+1):
                first = spectrum(L, ell, D('1.1'), D('.04'))
                second = spectrum(L, ell, D('1.1'), D('.02'))
                for k, (v1, v2) in enumerate(zip(first, second)):
                    slope = (v1/v2).ln()/D(2).ln()
                    largest_slope_error = max(
                        largest_slope_error, abs(slope-2*k))
                    assert abs(slope-2*k) < D('.08')
                    counts[k] += 2*ell+1
                    checks += 1
            assert counts == {k: (L-k+1)**2 for k in counts}
            checks += 1
            c = trig(D('1.1'))
            h = D(2)**(2*L)/D((2*L+1)*comb(2*L, L)**2)
            full_LL = moment(0, 2*L, c, semicircle)
            coefficient = trig(D('1.1'), True)**(2*L)*h/full_LL
            ratio = spectrum(L, 0, D('1.1'), D('.01'))[-1]/(
                coefficient*D('.01')**(2*L))
            assert abs(ratio-1) < D('.01')
            checks += 1
            thin.append({'L': L, 'coefficient': str(coefficient),
                         'ratio_at_0.01': str(ratio),
                         'max_slope_error': str(largest_slope_error)})

        for L in range(1, 8):
            p0 = [D(1)/D(n+1)**2 for n in range(L+1)]
            m0, m1 = mixing(L, D('.9')), mixing(L, D('.90001'))
            cov = [sum(a*b for a, b in zip(row, p0)) for row in m0]
            p1 = upper_solve(m1, cov)
            assert min(p1) > 0
            residual = max(abs(sum(a*b for a, b in zip(row, p1))-c)
                           for row, c in zip(m1, cov))
            max_covariance_residual = max(max_covariance_residual, residual)
            assert residual < D('1e-125')
            checks += 1

        # Closed-form kappa=2 curvature reconstruction and
        # independent top-mode ratio.
        for L in (1, 2, 4, 7):
            R, d1, d2 = D(10), D(3), D(6)
            g = (trig(d2/R, True)/trig(d1/R, True))**L
            assert abs(g-radial(L, L, d2/R)/radial(L, L, d1/R)) < D('1e-125')
            eta = g**(D(1)/D(L))
            # Invert cos by bisection; no trigonometric inverse dependency.
            lo, hi = D(0), PI/2
            for _ in range(450):
                mid = (lo+hi)/2
                if 2*trig(mid) > eta:
                    lo = mid
                else:
                    hi = mid
            recovered = d1/((lo+hi)/2)
            assert abs(recovered/R-1) < D('1e-100')
            curvature_rows.append(
                {'L': L, 'g': str(g), 'R_recovered': str(recovered)})
            checks += 1

        # Strict monotonicity, including the branch near the antipode.
        for kap in (D('1.2'), D(2), D('3.7')):
            for j in range(1, 100):
                x = PI*D(j)/(100*kap)
                derivative = (kap*trig(kap*x)/trig(kap*x, True)
                              - trig(x)/trig(x, True))
                assert derivative < 0
                checks += 1

        # Exact fractional Fisher Schur complement for deterministic
        # nuisance amplitudes.
        g, v1, v2, norm2 = Q(7, 5), Q(2), Q(3), Q(11)
        direct = norm2/v2-(g*g*norm2/v2**2)/(1/v1+g*g/v2)
        assert direct == norm2/(v2+g*g*v1)
        checks += 1

        report = {'status': 'PASS', 'checks': checks, 'decimal_digits': 160,
                  'rank_table': rank_rows, 'thin_shell': thin,
                  'max_isometry_residual': str(max_isometry),
                  'max_covariance_residual': str(max_covariance_residual),
                  'two_shell_reconstruction': curvature_rows}
    Path('shell_verification.json').write_text(json.dumps(report, indent=2))
    print(json.dumps(report, indent=2))


if __name__ == '__main__':
    verify()
\end{lstlisting}

\begin{thebibliography}{99}
\bibitem{P09} A. Preece and B. Batenin,
\emph{Conditional Reconstruction of a Quasi-Local Observable Domain
within Compact $S^3$ Geometry}, supplied unpublished manuscript,
file dated June 6, 2026.
\bibitem{P10} A. Preece and B. Batenin,
\emph{Conditional Reconstruction of a Quasi-Local Observable Domain
within Compact $S^3$ Geometry}, supplied unpublished manuscript,
file dated June 10, 2026 (internal date June 9, version 29).
\bibitem{CapDraft} A. Preece and B. Batenin,
\emph{Observability hierarchy for band-limited scalar fields on closed
spatial slices}, unpublished working manuscript, September 9, 2026.
\bibitem{Castell} W. zu Castell, N. La\'in Fern\'andez, and Y. Xu,
\emph{Polynomial Interpolation on the Unit Sphere II},
\href{https://doi.org/10.1007/s10444-005-7510-5}
{Adv. Comput. Math. \textbf{26}, 155 (2007)};
\href{https://arxiv.org/abs/math/0407448}{arXiv:math/0407448}.
\bibitem{DLMFQuad} NIST Digital Library of Mathematical Functions,
\emph{Quadrature}, Sec.~3.5,
\url{https://dlmf.nist.gov/3.5}, accessed September 10, 2026.
\bibitem{Simons} F. J. Simons, F. A. Dahlen, and M. A. Wieczorek,
\emph{Spatiospectral concentration on a sphere},
\href{https://doi.org/10.1137/S0036144504445765}
{SIAM Rev. \textbf{48}, 504 (2006)}.
\bibitem{Kinney} W. H. Kinney,
\emph{How to fool cosmic microwave background parameter estimation},
\href{https://doi.org/10.1103/PhysRevD.63.043001}
{Phys. Rev. D \textbf{63}, 043001 (2001)}.
\bibitem{Samandar} J. Noltmann \emph{et al.} (COMPACT Collaboration),
\emph{Cosmic topology. Part IIc. Detectability with non-standard primordial
power spectrum},
\href{https://arxiv.org/abs/2602.15527}
{arXiv:2602.15527 [astro-ph.CO] (2026)}.
\bibitem{Abramo} L. R. Abramo and K. E. Leonard,
\emph{Why multi-tracer surveys beat cosmic variance},
\href{https://doi.org/10.1093/mnras/stt465}
{Mon. Not. R. Astron. Soc. \textbf{432}, 318 (2013)};
\href{https://arxiv.org/abs/1302.5444}{arXiv:1302.5444 [astro-ph.CO]}.
\bibitem{DLMFPoly} NIST Digital Library of Mathematical Functions,
\emph{Definitions of classical orthogonal polynomials}, Sec.~18.3,
\url{https://dlmf.nist.gov/18.3}, accessed September 10, 2026.
\end{thebibliography}
\end{document}
